\scanpage{60}
% manuscript line 1
\deleted{\unsure{en groupant en produit, en écrivant} $u_1\otimes\cdots\otimes u$}
(car on a $u_1\otimes\cdots\otimes u_n=U_1\circ\cdots\circ U_n$, où $U_i$ est le
\deleted{$U_i=1\otimes\cdots u_i\otimes\cdots\otimes1$}
produit tensoriel de $n$ opérateurs, tous égaux à l'identité
sauf le $i$-ième, qui est égal à $u_i$. \deleted{Mais \unsure{identifiant les}}

% manuscript line 6
\deleted{Supposons donc $u\otimes1\otimes\cdots\otimes1$, où $A\in L(E_1,F_1)$,
identifiant $E_i$ à un espace $L^2(A_i)$, \unsure{avec} $F_i$ à un espace
$L^2(B_i)$, alors $\overset{(2)}{\bigotimes}_{1\leq i\leq n}E_i$ est identifié à l'espace $L^2(A_1\times\cdots\times A_n)$,
et $F_1\overset{(2)}{\otimes}E_{\unsure{2|1}}\overset{(2)}{\otimes}\cdots\overset{(2)}{\otimes}E_n$ est identifié à l'espace $L^2(A_1$}

% manuscript line 10
En vertu de l'associativité et de la commutativité du
produit tensoriel hilbertien, on \deleted{\unsure{peut}} est ramené alors au cas de
deux facteurs $E_1$ et $E_2$, $F_1$ et $F_2=E_2$, et où $A_1\in L(E_1,E_2)$, $A_2=\underline1$,
Mais alors, identifiant $E_1\overset{(2)}{\otimes}E_2$ à $L^{(2)}(\bar E_1',E_2)$, $F_1\overset{(2)}{\otimes}F_2$ à
$L^{(2)}(\bar F_1',E_2)$, \deleted{on est ramené} on est ramené
à prouver que pour $u\in L^2(\bar E_1',E_2)$, et $u$ de rang fini,
on a $\|u\,{}^tA_1\|_2\leq\|u\|_2\,\|{}^tA_1\|$, ce qui s'est en effet
avéré par la \deleted{formule de} deuxième formule (17) du N° 1,
aux notations près. Par suite, \deleted{l'indication $u_1\otimes\unsure{\cdots\otimes}$} \inserted{en vertu de (3),}
donc l'opération $\bigotimes_{1\leq i\leq n}u_i$ \deleted{\unsure{continue}} se \deleted{\unsure{tran}sf} prolonge
par continuité en une opération de $\overset{(2)}{\bigotimes}_{1\leq i\leq n}E_i$ dans
$\overset{(2)}{\bigotimes}_{1\leq i\leq n}F_i$, notée encore de la $=$ façon. \deleted{\unsure{Reprenons les}}

% manuscript line 22
\deleted{\unsure{Soient donc donnés} des espaces de Hilbert $E_1,\ldots,E_n$.}
\inserted{Soit maintenant} $E$ un espace de Hilbert, considérons
$\overset{n}{\otimes}\vphantom{\otimes}^{(2)}E$, \deleted{Posons \unsure{encore}} Posons

% manuscript line 25
\[
a_n=\frac1{n!}\sum_{\sigma\in\mathfrak S_n}\varepsilon_\sigma\,\sigma \tag*{(3)}
\]

% manuscript line 26
où la somme du deuxième mb est considérée comme
somme d'opérateurs de symétrie dans $\overset{n}{\otimes}\vphantom{\otimes}^{(2)}E$ ($\varepsilon_\sigma$ signature
de la permutation $\sigma$). \deleted{\unsure{un proj}} On reconnait, au facteur $\frac1{n!}$

% manuscript line 29
% archive sheet number: 61
